\section{Termékvízió}

A WatchDB termékvíziója egy olyan webalkalmazás létrehozása, amely egyszerre működik személyes filmnaplóként és később közösségi filmfelfedező platformként. A projekt alapgondolata, hogy a felhasználó egyetlen, platformfüggetlen helyen tudja rögzíteni, milyen filmeket nézett meg, miket szeretne később megnézni, és hogyan értékelte a már látott filmeket. Ez a megközelítés elsősorban a filmnézési előzmények rendszerezésére és a személyes filmnapló kialakítására épül.

A termék egyik célja a személyes használat. Ebben a felhasználó filmeket kereshet, részletes információkat tekinthet meg róluk, hozzáadhatja őket a watchlisthez, illetve megnézettként jelölheti és értékelheti őket. Ez a funkciókör adja a termék alapértékét, mert akkor is hasznos, ha a felhasználó egyedül használja az alkalmazást. A cél az, hogy a filmek naplózása gyors és egyszerű legyen, ne igényeljen hosszú szöveges értékeléseket, és ne legyen bonyolultabb annál, mint amit egy hétköznapi filmnéző szívesen használna.

A termék második, hosszabb távú célja a közösségi használat. A PRD alapján a WatchDB nemcsak azt tenné lehetővé, hogy a felhasználó saját filmjeit nyilvántartsa, hanem azt is, hogy lássa, barátai vagy ismerősei milyen filmeket néztek meg, hogyan értékelték azokat, illetve milyen filmeket tettek a saját listáikra. Ez a közösségi dimenzió a személyes filmnaplót ajánlási és felfedezési felületté bővítené: a felhasználó nem algoritmikus ajánlórendszeren keresztül, hanem ismerőseinek aktivitása alapján találhatna új filmeket.

Fontos különbséget tenni a teljes termékvízió és az elsőként megvalósítandó MVP között. A teljes vízió része lehet a publikus profil, a követési rendszer, az aktivitási feed, a reakciók, a statisztikai nézetek vagy akár egy fejlettebb ajánlórendszer is. Ezek azonban nem feltétlenül szükségesek ahhoz, hogy az alkalmazás első verziója értéket adjon. Az MVP ezért a termék személyes filmnapló funkcióira koncentrál: a keresésre, a filmrészletek megjelenítésére, a watchlist kezelésére, a megnézett filmek naplózására és az értékelésre.